% !TITLE 长相思·其一
% !AUTHOR 李白
% !DYNASTY 唐
% !GENRE 乐府诗
% !ORDER 14

\begin{poem}
长相思，在长安。
络纬\textsuperscript{1}秋啼金井阑\textsuperscript{2}，微霜凄凄簟色寒\textsuperscript{3}。
孤灯不明思欲绝，卷帷望月空长叹。
美人如花隔云端。
上有青冥\textsuperscript{4}之长天，下有渌水\textsuperscript{5}之波澜。
天长路远魂飞苦，梦魂不到关山难\textsuperscript{6}。
长相思，摧心肝。
\end{poem}

\begin{annotations}
\notetext{1}{络纬：纺织娘，一种秋虫，叫声像织布声。秋天的纺织娘在井栏边啼叫。}
\notetext{2}{金井阑：井边的栏杆。“金”是修饰词，形容栏杆的精美。}
\notetext{3}{簟色寒：竹席的颜色透着寒意。簟（diàn），竹席。微霜和簟寒一起写秋夜的清冷。}
\notetext{4}{青冥：青天、高空。冥（míng），深远。}
\notetext{5}{渌水：清澈的水。渌（lù），水清。}
\notetext{6}{关山难：关隘和山岭的阻隔，难以跨越。路太远，连梦魂都飞不过去。}
\end{annotations}

\begin{appreciation}
《长相思》是李白用乐府旧题写的组诗，这是第一首，写一个人深夜里想念远方的人。但细读你会发现，诗里的“美人”，恐怕不只是一个美人。

长相思，在长安。开头五个字，把相思和一座城绑在一起。李白在长安待过两年多，被赐金放还，从此再没回去。所以“在长安”三个字，是说给一个回不去的地方听的。接下来全是秋夜的细节：络纬秋啼金井阑，微霜凄凄簟色寒。络纬就是纺织娘，秋天叫起来像织布声；簟是竹席。秋虫在井栏边叫，薄霜落在竹席上——秋夜冷，一个人的秋夜更冷。孤灯不明思欲绝，卷帷望月空长叹。灯是孤灯，人是独坐；月亮在古诗里是相思的信使，可这封信没有收件人，只能空长叹。

美人如花隔云端。这是全诗最疼的一句。花，是眼前够得着的；云端，是永远够不着的。可望而不可即——读《蒹葭》时我们说过的词，伊人“在水一方”，李白把水换成了天，距离反而更远。上有青冥之长天，下有渌水之波澜。青冥是青天，渌水是清波。天高水阔，世界大得没边，却容不下两个人相见。天长路远魂飞苦，梦魂不到关山难。想得狠了连梦都做，可梦魂飞不过关山——路太远，连梦都会累。

结尾五个字：长相思，摧心肝。比开头重了一倍。

把这首诗和《行行重行行》放在一起读很有意思。古诗十九首里的女子也守着相思，守到最后还能说一句“努力加餐饭”，洒脱得很。李白连这种洒脱都做不到，梦到不了，话传不出，只能摧心肝。一个是忍住，一个是忍不住，都是真的相思。

再说那个“美人”。屈原以后，中国诗人有个老传统，用美人来写理想——还记得《离骚》里的香草美人吗？李白的“美人如花隔云端”，隔着云的恐怕还有他回不去的长安、放不下的抱负。长安和理想，都像花一样好看，都像云一样够不着。

这首诗哥哥每次读到“摧心肝”都要停一下。等你大了，也许会有想见却见不到的人，哥哥只希望那一天，你想念的人离你不算太远，想见就能见着——那就比李白强多了。
\end{appreciation}
